\section{Conclusions}
\label{sec:conclusions}

This work reframes SVGD as a numerically split interacting-particle flow in
which attraction to high-density regions and diversity-preserving repulsion are
treated as distinct computational components. From this viewpoint, multirate
integration is not only an efficiency heuristic but a practical mechanism for
improving stability in posterior geometries where these components evolve at
different effective scales.

Across a broad benchmark suite, the main empirical pattern is that multirate
formulations make SVGD behavior more reliable under difficult geometry,
especially in high-dimensional anisotropic settings and hierarchical models.
Fixed multirate schedules provide a strong robustness baseline, while adaptive
drift resolution offers additional flexibility when local stiffness changes
along the trajectory. At the same time, results indicate that no single variant
dominates every metric and every regime, reinforcing the importance of
evaluating quality and cost jointly rather than through a single summary
statistic.

Several limitations remain. Current adaptivity is centered on drift resolution,
while repulsion frequency is still controlled by simple schedules; ESS remains a
univariate proxy; and dense kernel interactions continue to constrain scaling at
large particle counts. These observations suggest clear next steps: tighter
adaptive controllers that coordinate drift and repulsion updates, approximate or
structured kernel computations for large-$N$ regimes, and stronger theoretical
links between multirate stability analysis and particle-based Bayesian
inference guarantees. Together, these directions can move multirate SVGD from a
robust empirical strategy toward a more principled and scalable foundation for
modern Bayesian computation.
